\scanpage{35}
$\operatorname{Aff}_A(G)$ opère \deleted{à gauche et à droite} sur $U_A(G)$ ; les formules sont \deleted{\unsure{de la forme}} analogues
\[
\langle g,f\mathbin{\lrcorner}X\rangle=\langle gf,X\rangle
\qquad\text{en particulier }\langle f,X\rangle=\varepsilon(f\mathbin{\lrcorner}X)
\]

Supposons maintenant que $A\subset B$, alors
\[
\operatorname{Aff}_B(G)=\operatorname{Aff}_A(G)\otimes_A B\supset\operatorname{Aff}_A(G)
\]
[on a une \uline{inclusion}, car $\operatorname{Aff}_A(G)$ plat sur $A$]
\[
U_B(G)=U_A(G)\otimes_A B\supset U_A(G)
\]
[raison analogue ; la première égalité est due au fait que $G$ est $A$-simple]

\uline{Proposition 1} Soit $X\in U_B(G)$. Conditions équivalentes
\[
\begin{array}{ll}
\text{i)}&X\in U_A(G)\\
\text{ii)}&X*\operatorname{Aff}_A(G)\subset\operatorname{Aff}_A(G)
\qquad\text{ii bis)}\ \operatorname{Aff}_A(G)*X\subset\operatorname{Aff}_A(G)\\
\text{iii)}&\langle X,\operatorname{Aff}_A(G)\rangle\subset A
\end{array}
\]
En effet,
\diagram{\begin{tikzcd}[column sep=small,row sep=small]
&\text{(ii bis)}\arrow[dr,Rightarrow]&\\
\text{(i)}\arrow[ur,Rightarrow]\arrow[r,Rightarrow]&\text{(ii)}\arrow[r,Rightarrow]&\text{(iii)}
\end{tikzcd}}
triviaux.

(iii) implique que \deleted{$\unsure{X\in U_A}$} \deleted{la forme} \inserted{$X$ définit} une forme $A$-linéaire sur $\operatorname{Aff}_A(G)$ nulle sur une puissance de l'idéal d'augmentation, donc un élément $X_0$ de $U_A(G)$, qui \unsure{n'est} évidemment \unsure{autre que} $X$
\margin{34}
